\documentclass[12pt]{amsart}


\usepackage{amssymb}
\usepackage{amsthm}
\usepackage{amsmath}
\usepackage{amsxtra}
\usepackage{mathrsfs}
\usepackage{enumerate}
\usepackage{colonequals}

\setlength{\hfuzz}{4pt}

\newenvironment{problab}[1]
{\noindent\textbf{Problem #1}.}
{\vskip 6pt}

\newenvironment{exolab}[1]
{\noindent\textbf{Exercise #1}.}
{\vskip 6pt}


\theoremstyle{remark}
\newtheorem*{soln}{Solution}

\newcommand{\C}{\mathbb C}
\newcommand{\D}{\mathfrak D}
\renewcommand{\H}{\mathfrak H}
\newcommand{\F}{\mathbb F}
\renewcommand{\P}{\mathbb P}
\newcommand{\PP}{\mathbb P}
\newcommand{\Q}{\mathbb Q}
\newcommand{\R}{\mathbb R}
\renewcommand{\SS}{\mathbb S}
\newcommand{\Z}{\mathbb Z}

\renewcommand{\Im}{\operatorname{Im}}
\renewcommand{\Re}{\operatorname{Re}}

\newcommand{\eps}{\epsilon}

\usepackage{fullpage}

\DeclareMathOperator{\Aut}{Aut}
\DeclareMathOperator{\Hom}{Hom}
\DeclareMathOperator{\ord}{ord}

\usepackage{mathpazo}

\begin{document}

\title{Belyi Maps and Dessins d'Enfants \\ Homework \#3}
\date{\today}

%\thanks{\emph{Date:} Monday, 11 February 2013}

\maketitle

\begin{exolab}{3.1}
  Let $C: X^3 + Y^3 = Z^3$ and define
  \begin{align*}
    f: C &\to \P^1\\
    [X:Y:Z] &\mapsto [X:Z] \, .
  \end{align*}
  Determine the ramification points of $f$ and their multiplicities.
\end{exolab}

\begin{exolab}{3.2}
  \begin{enumerate}[(a)]
    \item 
      Let $f \in \mathcal{M}(X)$ be a nonconstant meromorphic function on a compact Riemann surface $X$. Show that $\sum_P \ord_P(f) = 0$. (\textit{Hint}: Use results about degree.)
    \item
      Given $0 \neq x \in \Z$ and a prime $p \in \Z$, define the \textsf{$p$-adic valuation}
      \begin{equation*}
        \ord_p(x) \colonequals \max\{v \in \Z : p^v \mid x\}
      \end{equation*}
      and extend the definition to $x = a/b \in \Q$ by
      \begin{equation*}
        \ord_p(x) = \ord_p(a/b) \colonequals \ord_p(a) - \ord_p(b) \, .
      \end{equation*}
      By convention, we define $\ord_p(0) = \infty$. Define the \textsf{$p$-adic absolute value} by
      \begin{equation*}
        |x|_p = p^{-\ord_p(x)}
      \end{equation*}
      for $x \in \Q$. Finally, define $|x|_\infty = |x|$, the usual absolute value. Show that
      \begin{equation*}
        \prod_{p \leq \infty} |x|_p = 1
      \end{equation*}
      where the product ranges over all primes $p \in \Z$ as well as $p = \infty$.
  \end{enumerate}
\end{exolab}


\begin{exolab}{3.3}
  Let $C$ be a hyperelliptic curve given by the Weierstrass equation
  \begin{equation*}
    Y^2 = f(X,Z) = \prod_{i=1}^{2g+2} (X - \alpha_i Z)
  \end{equation*}
  in $\PP(1,g+1,1)$, where the $\alpha_i \in \C$ are distinct.
  \begin{enumerate}[(a)]
    \item 
      Let $U_0$ and $U_2$ be the affine open subsets where $X \neq 0$ and $Z \neq 0$, respectively. Write $C \cap U_0$ and $C \cap U_2$ as affine curves given by (affine) Weierstrass equations.
    \item
      Determine the relationship between your local affine coordinates on $U_0$ and $U_2$. E.g., if $x,y$ are affine coordinates on $U_0$ and $z,w$ are affine coordinates on $U_2$, express $z$ and $w$ as rational functions in $x$ and $y$, valid on $U_0 \cap U_2$.
    \item
      Compare our definition of a hyperelliptic curve with the one given in Miranda on p. 60. Do you see how they are equivalent?
  \end{enumerate}
\end{exolab}

\end{document}
